\clearpage
\hypertarget{native-475-lower}{}
\begin{center}
\rule{0.55\linewidth}{0.4pt}\\[2pt]
Astronomische Nachrichten Band XXII. Nr. 526. S. 353. 354. 25. Januar 1845.\\[2pt]
\rule{0.55\linewidth}{0.4pt}
\end{center}

Am 11. verglich Herr \textsc{Petersen} den Kometen mit zwei Sternen aus der Hist. Cél., S. 237 und 295, für die er zur Beobachtungszeit die folgenden scheinbaren Positionen annahm

\begin{center}
\begin{tabular}{rr}
$19^{\rm u}\ 21^{\rm m}\ 2^{\rm s}85$ & $+42^\circ\ 1'\ 5''1$ \\
$22\ 20.95$ & $41\ 55\ 11.8$ \\
\end{tabular}
\end{center}
\noindent Er fand:

\begin{center}
Jan. 11.\quad $9^{\rm u}\ 16^{\rm m}\ 7^{\rm s}$ Altonaer m. Z.\quad $289^\circ\ 37'\ 35''6$\quad $+41^\circ\ 58'\ 40''7$.
\end{center}

Aus zwei Vergleichen mit denselben Sternen erhielt Herr Hofrat \textsc{Gauss}

\begin{center}
\begin{tabular}{lll}
Jan. 11. & $5^{\rm u}\ 35^{\rm m}\ 19^{\rm s}2$ Gött. m. Z. & $+41^\circ\ 54'\ 40''2$ \\
 & $5\ 39\ 29.0$ & Rektasz. $289^\circ\ 41'\ 47''6$ \\
\end{tabular}
\end{center}

Die Deklination beruht allein auf dem zweiten Stern, die Rektaszension auf dem Mittel aus beiden. Herr Hofrat \textsc{Gauss} wünscht, dass die Positionen der Sterne jetzt bei der unteren Kulmination bestimmt werden.

Zum Schluss dieses Artikels erhalte ich noch von Herrn Hofrat \textsc{Gauss} eine Göttinger Beobachtung, von Herrn \textsc{Rümker} die folgenden Elemente und von Herrn \textsc{Sievers} die Verbesserung seiner früheren Bahn. Herr \textsc{Rümker} hätte seine Elemente früher geliefert, wenn die Geschäfte der Navigationsschule, bei denen ihm mehr Hilfe zu wünschen wäre, nicht gerade fast seine ganze Zeit in Anspruch nähmen.

\begin{center}
\setlength{\tabcolsep}{3pt}
\begin{tabular}{lll}
1845 Jan. 12. & $5^{\rm u}\ 51^{\rm m}\ 11^{\rm s}0$ M. Z. Göttingen & Rektasz. $289^\circ\ 14'\ 45''5$ \\
 & $5\ 54\ 59.5$ & Dekl. $42\ 22\ 15.4$ N. \\
\end{tabular}
\end{center}

\Needspace{7\baselineskip}
\SourceLeaf{476}
Zum Vergleich dienten 3 Sterne. Ihre aus der Hist. Cél. abgeleiteten Positionen für den Beobachtungstag sind:

\begin{center}
\begin{tabular}{rr}
$288^\circ\ 54'\ 5''4$ & $+42^\circ\ 15'\ 49''7$ \\
$289\ 1\ 5.6$ & $42\ 23\ 59.7$ \\
$290\ 7\ 35.2$ & $42\ 29\ 48.2$ \\
\end{tabular}
\end{center}

Der letzte dieser Sterne ist auf \textsc{Hardings} Karte deutlich zu weit nördlich eingezeichnet.
\begin{flushright}\textsc{Schumacher}.\end{flushright}
